\subsection{The averaged cube estimate for the three-dimensional
$\U(1)$ Villain model: exact duality, dual form, and the missing
estimate}\label{rM3:sec}

Theorem~\ref{rD:maxwell} shows that the free photon fails
$\mathrm{BP}(n,\varpi)$ at any fixed $n$ as $L\to\infty$, and
Remark~\ref{rem:rSmissing} records that step (D) needs a cube constant
produced by a dynamically generated mass. A lattice gauge theory in
which confinement and exponential decay of correlations are proved at
every coupling, including weak coupling, is the three-dimensional $\U(1)$
Villain model (G\"opfert--Mack~\cite{GM82}); in the semiclassical picture
of Polyakov~\cite{Pol77} and Banks--Myerson--Kogut~\cite{BMK77} its photon
acquires the Debye mass $m_D\sim\exp(-c\beta)$ of a dilute monopole gas.
This subsection sets up the exact duality of that model in the form needed
for the cube dynamics, identifies what $\mathrm{BP}(n,\varpi)$ becomes on
the dual side, proves it for the Debye--H\"uckel Gaussian approximation,
and isolates what is missing for the Villain measure itself. Nothing here
proves $\mathrm{BP}(n,\varpi)$ for the Villain model. The results are: the
exact decomposition of the field strength into an independent massless
Gaussian and a lattice-valued monopole field
(Proposition~\ref{rM3:lawF}); the exact dual form of the condition on the
plaquette algebra (Theorem~\ref{rM3:dualform}); the law of the interior
field given the exterior (Proposition~\ref{rM3:interior}), whose exact
part is a Gaussian independent of the exterior; the inequality of
$\mathrm{BP}$ with a $\beta$-independent constant for the odd
single-plaquette functions (Proposition~\ref{rM3:plaq}); a two-sided bound
$c\min(1,s^2)\le\varpi\le C\min(1,s^2)$, $s=nm$, uniform in $L$ and
$\beta$, for the Debye--H\"uckel Gaussian field of mass $m$
(Theorem~\ref{rM3:DH}); the
necessary condition $\varpi\le e^2q_nn/\xi$ in terms of the correlation
length (Proposition~\ref{rM3:xi}); and the precise statement of the open
estimates (Sec.~\ref{rM3:missing}).

\subsubsection{Setting}\label{rM3:setting}

The lattice is the torus $\Lambda_L=(\Z/L\Z)^3$, $L$ even, $L\ge2n+4$,
$n\ge2$; $E$, $P$, $C$ are its links, plaquettes and (unit) cubes, with
$\lvert E\rvert=\lvert P\rvert=3L^3$, $\lvert C\rvert=L^3$. Conventions
are those of Sec.~\ref{rD:sec} with $4$ replaced by $3$: a plaquette
$(y,\mu\nu)$, $\mu<\nu$, has links $(y,\mu),(y+e_\mu,\nu),(y+e_\nu,\mu),
(y,\nu)$ and
\[
(d\theta)_{(y,\mu\nu)}=\theta_{(y,\mu)}+\theta_{(y+e_\mu,\nu)}
-\theta_{(y+e_\nu,\mu)}-\theta_{(y,\nu)};
\]
a cube $c_z$ has corner $z$ and $(dF)_z=\sum_{\mu<\nu}\pm(F_{(z+e_\rho,
\mu\nu)}-F_{(z,\mu\nu)})$ with the sign $+$ for $(\mu\nu,\rho)=(23,1),
(12,3)$ and $-$ for $(13,2)$; $\delta=d^{\mathsf T}$ for the standard inner
products on $\R^E,\R^P,\R^C$. Then $dd=0$, $\delta\delta=0$, and
$-\Delta_3=d\delta$ on $\R^C$ is the (positive) lattice Laplacian of the
dual lattice with kernel the constants. For $V\subset E$, $P(V)$ is the set
of plaquettes containing a link of $V$; $V_x$ is the set of links with
base point in $x+[0,n)^3$, so $\lvert P(V_x)\rvert=3n^2(n+2)$; $X=
\Lambda_L$; $q_n=(2+1/n)^3\in(8,\frac{125}8]$. The gauge group is
$\U(1)^{\Lambda_L}$, $\mathcal A\subset L^2(\mu)$ the gauge-invariant
functions, $\mathcal P_x$ the conditional expectation given the links
outside $V_x$, and $\mathrm{BP}(n,\varpi)$ is
\[
\varpi\Var_\mu(f)\le\mathcal E_n(f):=\frac1{n^3}\sum_{x\in X}
\mu\bigl((f-\mathcal P_xf)^2\bigr)
\]
for every real $f\in\mathcal A$.

\begin{remark}[dimension]\label{rM3:dim}
Lemma~\ref{lem:avgspeed} and Theorem~\ref{thm:avgdecay} use, besides
the facts that $\mu$ has a positive density and that each $\mathcal P_x$
is an orthogonal projection preserving $\mathcal A$, only the facts that
each link lies in exactly $n^d$ of the sets $V_x$, that
$\lvert X\rvert\ge4n^d$ (here $\lvert X\rvert=L^3\ge8n^3$), and that
$y\rightsquigarrow z$ forces $\lvert z-y\rvert_\infty\le n$, so that at
most $(2n+1)^d=q_nn^d$ points feed one cube. Their proofs hold with
$d=3$ after one numerical change: where $q_n>16$ is used, $q_n>8$ is
used instead, which replaces $64e^2$ by $16e^2$ in the prefactor of
Theorem~\ref{thm:avgdecay} and leaves the inequalities
$\varpi/(e^2q_n)<1-e^{-2}$ and $3\varpi/(e^2q_n)<1$ valid. So
$\varpi\le4$ whenever $\mathrm{BP}(n,\varpi)$ holds (indeed $\varpi\le
1+2/n$, since a plaquette lies in $P(V_x)$ for $n^2(n+2)$ points $x$), and
$\mathrm{BP}(n,\varpi)$ gives covariance decay of local gauge-invariant
functions at rate $\varpi/(e^2q_nn)$ per lattice unit. The passage to a
transfer-matrix gap (Corollary~\ref{cor:avggap}) needs in addition the
reflection positivity of the three-dimensional Villain measure and the
three-dimensional form of the transfer-matrix construction; they are not
discussed here and are not used below.
\end{remark}

The Villain weight is $\varphi_\beta(t)=\sum_{k\in\Z}e^{-\beta(t+2\pi k)^2
/2}$, $t\in(-\pi,\pi]$, as in Proposition~\ref{prop:coulomb}, and
\[
\mu_\beta(d\theta)=\frac1Z\prod_{p\in P}\varphi_\beta\bigl((d\theta)_p
\bigr)\prod_{e\in E}\frac{d\theta_e}{2\pi},\qquad\theta\in(-\pi,\pi]^E .
\]
Write $P_\beta$ for the law of the field strength $F=d\theta+2\pi k$
introduced below; $\E$ and $\Var$ without subscript refer to $\mu_\beta$
or to the joint measure $\tilde\mu_\beta$, as the context shows.

\subsubsection{The joint measure and the law of the field strength}
\label{rM3:duality}

Let $\tilde\mu_\beta$ be the probability measure on $(-\pi,\pi]^E\times
\Z^P$ with density proportional to $\prod_{p}e^{-\beta((d\theta)_p+2\pi
k_p)^2/2}$ with respect to $\prod_e d\theta_e/2\pi\otimes$ counting
measure; summing over $k$ gives $\mu_\beta$ as the $\theta$-marginal. Put
$F:=d\theta+2\pi k\in\R^P$. Then $dF=2\pi\,dk$, $dk=:\mathsf n\in\Z^C$ is
the monopole density, $\sum_z\mathsf n_z=0$, and the periods of $F$ over
closed $2$-chains lie in $2\pi\Z$.

Let $\mathcal B=d\R^E\subset\R^P$ (the exact $2$-cochains), $\Pi$ the
orthogonal projection onto $\mathcal B$, $\mathcal K_\R=\ker d\subset\R^E$
and $\mathcal H=\ker d\cap\ker\delta\subset\R^P$ (harmonic $2$-cochains,
$\dim\mathcal H=3$). The Hodge decomposition is $\R^P=\mathcal B\oplus
\delta\R^C\oplus\mathcal H$ and $\mathcal B^\perp=\delta\R^C\oplus\mathcal
H=\ker\delta$.

\begin{lemma}[sector lattices]\label{rM3:lattices}
\textup{(i)} $\mathcal K_\Z:=\mathcal K_\R\cap\Z^E$ is a lattice of full
rank in $\mathcal K_\R$. \textup{(ii)} $\Lambda_{\mathcal B}:=\mathcal B
\cap\Z^P$ is a lattice of full rank in $\mathcal B$ containing $d\Z^E$ with
finite index. \textup{(iii)} $\mathcal M:=2\pi(1-\Pi)\Z^P$ is a lattice
(discrete subgroup of full rank) in $\mathcal B^\perp$.
\end{lemma}

\begin{proof}
(i) and (ii): $d$ has integer entries, so $\mathcal K_\R$ and $\mathcal B$
are rational subspaces (spanned by vectors with rational coordinates:
the kernel of an integer matrix has a rational basis, and $\mathcal B$ is
spanned by the integer columns of $d$). A rational subspace $W\subset\R^N$
of dimension $r$ meets $\Z^N$ in a lattice of rank $r$: clearing
denominators of a rational basis gives $r$ independent integer vectors
of $W$, and $W\cap\Z^N$ is discrete. $d\Z^E\subset\Lambda_{\mathcal B}$ has
rank $\dim\mathcal B$, so the index is finite. (iii) Let $B$ be the integer matrix whose columns are a basis of
$\mathcal B$ chosen among the columns of $d$; then $\Pi=B(B^{\mathsf T}B)
^{-1}B^{\mathsf T}$ is a rational matrix, so $(1-\Pi)\Z^P\subset N^{-1}
\Z^P$ for a common denominator $N$. Hence $(1-\Pi)\Z^P$ is a discrete,
finitely generated subgroup of $\mathcal B^\perp$, and it spans
$\mathcal B^\perp$ because $(1-\Pi)$ maps $\R^P$ onto $\mathcal B^\perp$
and $\Z^P$ spans $\R^P$.
\end{proof}

\begin{proposition}[law of the field strength]\label{rM3:lawF}
Under $\tilde\mu_\beta$ put $G:=\Pi F\in\mathcal B$ and $M:=(1-\Pi)F\in
\mathcal B^\perp$. Then
\begin{enumerate}[label=\textup{(\roman*)},leftmargin=*]
\item $M=2\pi(1-\Pi)k\in\mathcal M$, and $G=\Pi F=d\theta+2\pi\Pi k$;
\item $G$ and $M$ are independent; $G$ has the Gaussian law $\gamma_\beta$
on $\mathcal B$ with covariance $\beta^{-1}\Pi$ (the Maxwell field of
Theorem~\ref{rD:maxwell} in three dimensions), and
$P(M=m)=Z_{\mathcal M}^{-1}e^{-\beta\lVert m\rVert^2/2}$, $m\in\mathcal M$;
\item $F=G+M$, so $P_\beta=\gamma_\beta*\bigl(Z_{\mathcal M}^{-1}
\sum_{m\in\mathcal M}e^{-\beta\lVert m\rVert^2/2}\delta_m\bigr)$.
\end{enumerate}
\end{proposition}

\begin{proof}
(i) is immediate from $F=d\theta+2\pi k$ and $\Pi d\theta=d\theta$.
For (ii) let $\Phi\ge0$ be a Borel function on $\R^P$. By definition of
$\tilde\mu_\beta$,
\begin{align*}
Z\,\tilde\mu_\beta\bigl(\Phi(F)e^{\beta\lVert F\rVert^2/2}\bigr)
&=\int_{(-\pi,\pi]^E}\sum_{k\in\Z^P}\Phi(d\theta+2\pi k)\prod_e\frac{d\theta_e}
{2\pi}\\
&=\int_{\mathbb T^E}\bar\Phi(d\theta)\,d\theta,
\end{align*}
where $\mathbb T^E=(\R/2\pi\Z)^E$ carries normalised Haar measure and
$\bar\Phi(F+2\pi\Z^P)=\sum_k\Phi(F+2\pi k)$ is the periodisation of $\Phi$
to $\mathbb T^P$. The map $d:\mathbb T^E\to\mathbb T^P$ is a continuous
homomorphism, so the image of Haar measure is Haar measure on the compact
group $d(\mathbb T^E)=(\mathcal B+2\pi\Z^P)/2\pi\Z^P\cong\mathcal B/2\pi
\Lambda_{\mathcal B}$, a torus by Lemma~\ref{rM3:lattices}(ii), whose Haar
measure is $\lambda_{\mathcal B}$ restricted to a fundamental domain
$Q$ of $2\pi\Lambda_{\mathcal B}$ and divided by $\lvert Q\rvert$. For
$g\in\mathcal B$ group the terms of $\bar\Phi(g)=\sum_{k\in\Z^P}\Phi(g+2\pi
k)$ according to $m=2\pi(1-\Pi)k\in\mathcal M$: the $k$ with a given $m$
form a coset $k_0(m)+\Lambda_{\mathcal B}$, since $\Lambda_{\mathcal B}
=\Z^P\cap\ker(1-\Pi)$, and $2\pi k=2\pi\Pi k_0(m)+2\pi\ell+m$ for $\ell\in
\Lambda_{\mathcal B}$. Hence
\begin{align*}
\int_{\mathbb T^E}\bar\Phi(d\theta)\,d\theta&=\frac1{\lvert Q\rvert}
\sum_{m\in\mathcal M}\sum_{\ell\in\Lambda_{\mathcal B}}\int_Q\Phi\bigl(g+2\pi
\Pi k_0(m)\\
&\qquad\qquad+2\pi\ell+m\bigr)\lambda_{\mathcal B}(dg)\\
&=\frac1{\lvert Q\rvert}
\sum_{m\in\mathcal M}\int_{\mathcal B}\Phi(g+m)\,\lambda_{\mathcal B}(dg),
\end{align*}
because the translates $Q+2\pi\ell+2\pi\Pi k_0(m)$, $\ell\in\Lambda_{\mathcal
B}$, tile $\mathcal B$. Taking $\Phi(F)=\mathbf 1_A(\Pi F)\mathbf 1_{\{(1-\Pi)F=m
\}}e^{-\beta\lVert F\rVert^2/2}$ with $A\subset\mathcal B$ Borel and using
$\lVert g+m\rVert^2=\lVert g\rVert^2+\lVert m\rVert^2$ for $g\in\mathcal
B$, $m\in\mathcal B^\perp$,
\[
\tilde\mu_\beta(G\in A,\,M=m)=\frac{e^{-\beta\lVert m\rVert^2/2}}{Z\lvert Q
\rvert}\int_Ae^{-\beta\lVert g\rVert^2/2}\lambda_{\mathcal B}(dg),
\]
which factorises in $(A,m)$; this is (ii), and (iii) restates it.
\end{proof}

\begin{remark}[Coulomb form]\label{rM3:coulomb}
Each $m\in\mathcal M$ is $m=2\pi(\delta u+h)$ with $u\in\R^C$ of mean
zero, $h\in\mathcal H$, and $\mathsf n:=dk=\frac1{2\pi}dm=d\delta u=
-\Delta_3u$, so $u=(-\Delta_3)^{-1}\mathsf n$ and
\[
\lVert m\rVert^2=4\pi^2\bigl(\ip{\mathsf n}{(-\Delta_3)^{-1}\mathsf n}
+\lVert h\rVert^2\bigr)=:4\pi^2\bigl(\mathcal E(\mathsf n)+\lVert h
\rVert^2\bigr),
\]
the lattice Coulomb energy of the integer charge distribution $\mathsf n$
plus a harmonic part. The weight $e^{-\beta\lVert m\rVert^2/2}$ is
$e^{-2\pi^2\beta\mathcal E(\mathsf n)}e^{-2\pi^2\beta\lVert h\rVert^2}$:
this is the monopole Coulomb gas of \cite{BMK77,GM82}, with the
self-energy $2\pi^2\beta v_L(0)$ of an isolated unit charge,
$v_L=(-\Delta_3)^{-1}$ on mean-zero functions; as $L\to\infty$,
$v_L(0)$ tends to the value $v(0)=0.2527\ldots$ of the Green function
of $\Z^3$ at the origin (a standard fact, not used below). The harmonic
part is constrained by the integrality of $k$ jointly with $\mathsf n$,
so we do not claim independence of $\mathsf n$ and $h$; nothing below
uses it. In the Debye--H\"uckel approximation of a dilute gas, each
cube carries charge $\pm1$ with probability $z=e^{-2\pi^2\beta v(0)}$, so
$\mathsf n_z$ has variance $2z$, and the screening mass of the coupling
$4\pi^2\beta$ is
\[
m_D^2=4\pi^2\beta\cdot2z=8\pi^2\beta\,e^{-2\pi^2\beta v(0)}
\]
in lattice units~\cite{Pol77,BMK77}. This formula is heuristic and is
not used in any proof below; Definition~\ref{rM3:DHdef} turns the
Debye--H\"uckel approximation into a Gaussian measure with a free mass
parameter. The proof of Theorem~\ref{rD:maxwell} applies to $G$ in
three dimensions, with Lemma~\ref{rM3:layer} below (for $dF=0$) in place
of Lemmas~\ref{rD:layer} and~\ref{rD:walls}: it uses only $dG=0$ and
the covariance $\beta^{-1}\Pi$, and with $n$ layers of $n^2+2n$
plaquettes and $\operatorname{Var}(f)=L^3/(2\beta)$ it shows that the
component $G$ alone, with the exterior $\sigma$-algebras
$\sigma(G_p:p\notin P(V_x))$, fails $\mathrm{BP}(n,\varpi)$ unless
$\varpi\le8\pi^2n(n+2)/L^2$; this is the bound of
Theorem~\ref{rM3:DH}(iii) with $\mathsf m=0$. The
content of the duality is that the Villain field strength is exactly
this massless Gaussian plus an independent lattice-valued monopole field,
and that all screening is carried by the second summand.
\end{remark}

\subsubsection{Charge sectors and the plaquette algebra}\label{rM3:sectors}

For $c\in(-\pi,\pi]^3$ let $T_c$ shift every link of direction $\mu$ by
$c_\mu$: $(T_c\theta)_{(y,\mu)}=\theta_{(y,\mu)}+c_\mu$ mod $2\pi$. Then
$d(T_c\theta)=d\theta$, so $T_c$ preserves $\mu_\beta$ and the
plaquette variables; it maps gauge-invariant functions to gauge-invariant
functions, and it commutes with $\mathcal P_x$, since it maps the
conditional law of $\theta_{V_x}$ given the exterior to the conditional
law given the shifted exterior (the density $\prod_{p\in P(V_x)}\varphi_\beta
((d\theta)_p)$ is $T_c$-invariant). The torus $\mathbb T^3=\{T_c\}$ acts
unitarily on $L^2(\mu_\beta)$ and on $\mathcal A$; let $\mathcal A=
\bigoplus_{\kappa\in\Z^3}\mathcal A_\kappa$ be the decomposition into
isotypic components, $T_cf=e^{i\ip\kappa c}f$ on $\mathcal A_\kappa$. Each
$\mathcal A_\kappa$ is closed, $\mathcal P_x$-invariant, and the
components are mutually orthogonal. Let $\mathcal A_F$ be the closed
subspace of $L^2(\mu_\beta)$ of functions of the plaquette angles
$(e^{i(d\theta)_p})_{p\in P}$, that is, of $(e^{iF_p})_{p\in P}$.
Then $\mathcal A_F\subset\mathcal A_0$, and $\mathcal P_x\mathcal A_F
\subset\mathcal A_F$ by Corollary~\ref{rM3:exterior} below.

\begin{lemma}[splitting]\label{rM3:split}
Let $S\subset\mathcal A$ be a closed subspace containing the constants,
invariant under complex conjugation and under every $\mathcal P_x$,
$x\in X$. Then $\mathrm{BP}(n,\varpi)$
holds for $\mu_\beta$ if and only if $\varpi\Var(f)\le\mathcal E_n(f)$
for all real $f\in S$ and for all real $f\in S^\perp\cap\mathcal A$.
\end{lemma}

\begin{proof}
The direct implication is trivial. Conversely, let $f\in\mathcal A$ be
real and write $f=f_1+f_2$ with $f_1$ the orthogonal projection of $f$ onto
$S$ and $f_2=f-f_1\in S^\perp\cap\mathcal A$. Since $S$ is invariant
under conjugation, so is the projection onto it, and $f_1$, $f_2$ are real.
Then $\Var f=\Var f_1+\Var f_2$ since $f_2\perp S\ni1,f_1$, and
$\mathcal E_n(f)=\mathcal E_n(f_1)+\mathcal E_n(f_2)$ since
$f_1-\mathcal P_xf_1\in S$ and $f_2-\mathcal P_xf_2\in S^\perp$
($\mathcal P_x$ is self-adjoint and $S$-invariant, hence
$S^\perp$-invariant); the two inequalities add up.
\end{proof}

Applied to $S=\mathcal A_F$, the lemma reduces $\mathrm{BP}(n,\varpi)$ to
two statements: the inequality on functions of the field strength
(treated in Theorem~\ref{rM3:dualform}), and the inequality on
$\mathcal A_F^\perp\cap\mathcal A$, which contains the charged sectors
$\mathcal A_\kappa$, $\kappa\ne0$ (Polyakov loops winding around the
torus), and whatever part of $\mathcal A_0$ is not generated by the
plaquette angles. On a box with free boundary every gauge-invariant
function is a function of the plaquette angles (axial gauge fixing), and
on the torus one expects $\mathcal A_0=\mathcal A_F$; we do not prove
this and do not use it. The charged sectors are left open here (see
Remark~\ref{rM3:charged}).
\subsubsection{The dual block and the exterior data of a cube}
\label{rM3:block}

Fix $x\in X$ and use relative coordinates $z=y-x$, taken in
$(-L/2,L/2]^3$. The interior links are those with base point in
$[0,n)^3$; a plaquette $(y,\mu\nu)$ lies in $P(V_x)$ if and only if
$z_\rho\in[0,n-1]$ for the third direction $\rho$ and $(z_\mu,z_\nu)\in
[-1,n-1]^2\setminus\{(-1,-1)\}$ (its three base points $y,y+e_\mu,
y+e_\nu$ must include one in $[0,n)^3$). The cubes containing a plaquette
of $P(V_x)$ are the $c_{x+z}$ with
\[
z\in Y:=[-1,n-1]^3\setminus\{(-1,-1,-1)\},
\]
the dual block of the cube. Write $p(z,1)=(x+z+e_1,23)$, $p(z,2)=(x+z+e_2,
13)$, $p(z,3)=(x+z+e_3,12)$ for the plaquette common to $c_{x+z}$ and
$c_{x+z+e_\mu}$; then $p(z,\mu)\in P(V_x)$ exactly for the dual links
\begin{gather*}
(\dagger)\quad\{z,z+e_\mu\}:\ z_\mu\in[-1,n-2],\\
(z_\nu,z_\rho)\in[-1,n-1]^2\setminus\{(-1,-1)\},
\end{gather*}
and both ends of such a link lie in $Y$. For $u\in\R^C$ one has
$(\delta u)_{p(z,1)}=u_{z}-u_{z+e_1}$, $(\delta u)_{p(z,2)}=u_{z+e_2}-u_z$,
$(\delta u)_{p(z,3)}=u_z-u_{z+e_3}$ (writing $u_z$ for $u_{c_{x+z}}$).

\begin{lemma}[Poincar\'e lemma on the dual block]\label{rM3:poincare}
Let $m\in\R^P$ be supported in $P(V_x)$ with $(\delta m)_e=0$ for every
$e\in V_x$. Then there is $u\in\R^C$, supported on the cubes of $Y$, with
$(\delta u)_p=m_p$ for every $p\in P(V_x)$. If $m$ is integer valued, $u$
can be taken integer valued.
\end{lemma}

\begin{proof}
Put $a_1(z):=-m_{p(z,1)}$, $a_2(z):=m_{p(z,2)}$, $a_3(z):=-m_{p(z,3)}$ for
all $z$ (so $a_\mu(z)=0$ unless $\{z,z+e_\mu\}$ is a $(\dagger)$-link),
and for $\nu<\rho$
\[
\operatorname{curl}_{\nu\rho}(z):=a_\nu(z)+a_\rho(z+e_\nu)-a_\nu(z+e_\rho)
-a_\rho(z).
\]
Writing out $\delta$ on links with the conventions of
Sec.~\ref{rM3:setting} gives the identities
\begin{gather*}
(\delta m)_{(y,1)}=-\operatorname{curl}_{23}(z-e_2-e_3),\\
(\delta m)_{(y,2)}=\operatorname{curl}_{13}(z-e_1-e_3),\\
(\delta m)_{(y,3)}=-\operatorname{curl}_{12}(z-e_1-e_2),\qquad z=y-x .
\end{gather*}
For instance
\[
(\delta m)_{(y,1)}=m_{(y,12)}-m_{(y-e_2,12)}+m_{(y,13)}-m_{(y-e_3,13)},
\]
and with $z'=z-e_2-e_3$ the four terms are $-a_3(z'+e_2)$, $a_3(z')$,
$-a_2(z'+e_3)$ and $a_2(z')$. The hypothesis therefore reads
\[
(\star)\quad\operatorname{curl}_{\nu\rho}(z)=0\quad\text{if }
z_\mu\in[0,n-1],\ z_\nu,z_\rho\in[-1,n-2],
\]
where $\{\mu,\nu,\rho\}=\{1,2,3\}$.
We must find $u$ with $g_\mu(z):=u_{z+e_\mu}-u_z=a_\mu(z)$ on every
$(\dagger)$-link; the sign conventions above make this equivalent to
$(\delta u)_{p(z,\mu)}=m_{p(z,\mu)}$.

\emph{Step 1: the cube $Y_0=[0,n-1]^3$.} Put
\[
u_z:=\sum_{j<z_1}a_1(j,0,0)+\sum_{j<z_2}a_2(z_1,j,0)+\sum_{j<z_3}
a_3(z_1,z_2,j).
\]
Then $g_3(z)=a_3(z)$ directly. For $g_2$, the terms
\[
\sum_{j<z_3}[a_3(z_1,z_2+1,j)-a_3(z_1,z_2,j)]
\]
telescope by
$\operatorname{curl}_{23}(z_1,z_2,j)=0$ ($z_1\in[0,n-1]$, $z_2,j\in[0,n-2]$)
to $a_2(z)-a_2(z_1,z_2,0)$, so $g_2(z)=a_2(z)$. For $g_1$, $\operatorname{curl}_{12}(z_1,j,0)=0$ and
$\operatorname{curl}_{13}(z_1,z_2,j)=0$ give
\[
\begin{aligned}
&\sum_{j<z_2}[a_2(z_1+1,j,0)-a_2(z_1,j,0)]\\
&\qquad=a_1(z_1,z_2,0)-a_1(z_1,0,0),\\
&\sum_{j<z_3}[a_3(z_1+1,z_2,j)-a_3(z_1,z_2,j)]\\
&\qquad=a_1(z)-a_1(z_1,z_2,0);
\end{aligned}
\]
hence $g_1(z)=a_1(z)$. All curls
used are in the range $(\star)$.

\emph{Step 2 (faces).} For $z$ with $z_\mu=-1$ and the other two
coordinates in $[0,n-1]$ put $u_z:=u_{z+e_\mu}-a_\mu(z)$, so $g_\mu(z)
=a_\mu(z)$. For a link inside the face, say $\mu=1$ and $\{z,z+e_2\}$ with
$z_2\le n-2$: $g_2(z)=g_2(z+e_1)-a_1(z+e_2)+a_1(z)=a_2(z+e_1)-a_1(z+e_2)
+a_1(z)=a_2(z)+\operatorname{curl}_{12}(z)=a_2(z)$, since $z_3\in[0,n-1]$
and $z_1=-1,z_2\in[-1,n-2]$. The link $\{z,z+e_3\}$ uses
$\operatorname{curl}_{13}(z)=0$ ($z_2\in[0,n-1]$), and the faces $z_2=-1$,
$z_3=-1$ are treated in the same way with $\operatorname{curl}_{12}$,
$\operatorname{curl}_{23}$ and $\operatorname{curl}_{13}$,
$\operatorname{curl}_{23}$ respectively; in each case the coordinate that
must lie in $[0,n-1]$ is the one not involved in the face.

\emph{Step 3 (edges).} For $z$ with $z_\nu=z_\rho=-1$, $z_\mu\in[0,n-1]$,
put $u_z:=u_{z+e_\nu}-a_\nu(z)$ (the point $z+e_\nu$ lies on the face
$z_\rho=-1$). Then $g_\nu(z)=a_\nu(z)$, and
$g_\rho(z)=u_{z+e_\rho}-u_z=[u_{z+e_\nu+e_\rho}-a_\nu(z+e_\rho)]-[u_{z+e_\nu
+e_\rho}-a_\rho(z+e_\nu)-a_\nu(z)]=a_\rho(z)\pm\operatorname{curl}_{\nu\rho}
(z)=a_\rho(z)$, by $(\star)$ with $z_\mu\in[0,n-1]$. The links
$\{z,z+e_\mu\}$ along the edge are not $(\dagger)$-links, since
$(z_\nu,z_\rho)=(-1,-1)$.

\emph{Step 4.} Put $u_z:=0$ for $z=(-1,-1,-1)$ and for $z\notin Y$. The
links from the corner, the links along the edges, and the links with an
end outside $[-1,n-1]^3$ are not $(\dagger)$-links. Every $(\dagger)$-link
has $z\in Y$ with zero, one or two coordinates equal to $-1$, and it was
checked in Steps 1, 2, 3 respectively. Integrality is clear from the
construction.
\end{proof}

Let $d_x:\R^{V_x}\to\R^{P(V_x)}$ be the restriction of $d$ to functions
supported on $V_x$ (the interior coboundary), $D_x:=d_x\R^{V_x}$, $\Pi_x$ the
orthogonal projection of $\R^{P(V_x)}$ onto $D_x$, $K_x:=d_x(\mathbb T^{V_x})
\subset\mathbb T^{P(V_x)}$, a closed subgroup, and $\Gamma_x:=D_x+2\pi\Z^{P
(V_x)}$. The annihilator of $K_x$ in the dual group $\Z^{P(V_x)}$ of
$\mathbb T^{P(V_x)}$ is
\[
K_x^\perp=\{m\in\Z^{P(V_x)}:(\delta m)_e=0\ \forall e\in V_x\},
\]
because $\ip m{d_x\eta}=\ip{\delta m}\eta$ for $\eta\in\R^{V_x}$ and this
lies in $2\pi\Z$ for all $\eta$ exactly when $\delta m$ vanishes on $V_x$.
For $m\in K_x^\perp$ Lemma~\ref{rM3:poincare} gives $u^{(m)}\in\Z^C$ with
$m=\delta u^{(m)}$ on $P(V_x)$; put $m^{\prime}:=\delta u^{(m)}-m$, an
integer $2$-cochain supported outside $P(V_x)$.

\begin{lemma}[the coset of the exterior data]\label{rM3:coset}
For $\theta\in(-\pi,\pi]^E$ let $\theta^0$ agree with $\theta$ outside
$V_x$ and vanish on $V_x$, and let $b(\theta):=(d\theta^0)|_{P(V_x)}$. Then
for every $m\in K_x^\perp$
\begin{align*}
e^{i\ip m{b(\theta)}}&=\exp\Bigl(-i\sum_{p\notin P(V_x)}m'_p(d\theta)_p
\Bigr)\\
&=\exp\Bigl(-i\sum_{p\notin P(V_x)}m'_pF_p\Bigr),
\end{align*}
and the coset $b(\theta)+\Gamma_x$ is a continuous function of finitely
many of these quantities. Consequently $b(\theta)+\Gamma_x$ is a
measurable function of $(e^{iF_p})_{p\notin P(V_x)}$.
\end{lemma}

\begin{proof}
$\ip m{b(\theta)}=\ip m{d\theta^0}=\ip{\delta u^{(m)}}{d\theta^0}-\ip{m'}
{d\theta^0}=\ip{u^{(m)}}{dd\theta^0}-\sum_{p\notin P(V_x)}m'_p(d\theta)_p$,
since $m$ is supported in $P(V_x)$, $dd=0$ and $(d\theta^0)_p=(d\theta)_p$
for $p\notin P(V_x)$; the second equality uses $F_p-(d\theta)_p\in2\pi\Z$ and
$m'_p\in\Z$. The group $\mathbb T^{P(V_x)}/K_x$ is compact abelian with dual
group $K_x^\perp$ (Pontryagin duality), so its characters separate points;
$K_x^\perp$ is a subgroup of $\Z^{P(V_x)}$, hence finitely generated by
some $m_1,\dots,m_r$, and $\bar b\mapsto(e^{i\ip{m_j}{\bar b}})_{j\le r}$ is a
continuous injective homomorphism of a compact group into $\mathbb T^r$,
hence a homeomorphism onto its image. The coset $b+\Gamma_x$ is the class
of $b$ in $\mathbb T^{P(V_x)}/K_x$ when $\Gamma_x/2\pi\Z^{P(V_x)}=K_x$, which
holds since $K_x=d_x(\mathbb T^{V_x})=(D_x+2\pi\Z^{P(V_x)})/2\pi\Z^{P(V_x)}$.
\end{proof}

\begin{corollary}[the exterior data are the exterior field strength]
\label{rM3:exterior}
Let $f\in\mathcal A_F$. Then $\mathcal P_xf$ is a function of
$(e^{iF_p})_{p\notin P(V_x)}$, and under $\tilde\mu_\beta$
\begin{align*}
\mathcal P_xf&=\E\bigl[f\bigm|\theta_{E\setminus V_x},k_{P\setminus P(V_x)}
\bigr]=\E\bigl[f\bigm|F_{P\setminus P(V_x)}\bigr]\\
&=\E\bigl[f\bigm|(e^{iF_p})_{p\notin P(V_x)}\bigr]\quad\text{a.s.}
\end{align*}
In particular $\mathcal P_x\mathcal A_F\subset\mathcal A_F$.
\end{corollary}

\begin{proof}
Write $f=\Phi\bigl((e^{i(d\theta)_p})_{p\in P}\bigr)$. For $p\in P(V_x)$,
$(d\theta)_p=b(\theta)_p+(d_x\theta_{V_x})_p$, and for $p\notin P(V_x)$,
$(d\theta)_p$ depends only on $\theta_{E\setminus V_x}$. The conditional
law of $\theta_{V_x}$ given the exterior has density proportional to
$\prod_{p\in P(V_x)}\varphi_\beta(b(\theta)_p+(d_x\theta_{V_x})_p)$, so
$\mathcal P_xf=\Psi\bigl(b(\theta),(e^{i(d\theta)_p})_{p\notin P(V_x)}\bigr)$
with
\begin{gather*}
\Psi(b,w)=\frac{\int\Phi\bigl(e^{i(b+d_x\eta)},w\bigr)\Upsilon_b(\eta)\,d\eta}
{\int\Upsilon_b(\eta)\,d\eta},\\
\Upsilon_b(\eta)=\prod_{p\in P(V_x)}\varphi_\beta\bigl(b_p+(d_x\eta)_p\bigr),
\end{gather*}
the integrals over $\eta\in(-\pi,\pi]^{V_x}$.
The substitution $\eta\mapsto\eta-\eta_0$ shows $\Psi(b+d_x\eta_0,w)=
\Psi(b,w)$, and $\Psi(b+2\pi j,w)=\Psi(b,w)$ for $j\in\Z^{P(V_x)}$ since
$\Phi$ and $\varphi_\beta$ are $2\pi$-periodic. So $\Psi(\cdot,w)$ is a
function of the coset $b+\Gamma_x$, and the first claim follows from
Lemma~\ref{rM3:coset} ($e^{i(d\theta)_p}=e^{iF_p}$).

Under $\tilde\mu_\beta$ the conditional density of $(\theta_{V_x},k_{P(V_x)})$
given $(\theta_{E\setminus V_x},k_{P\setminus P(V_x)})$ is proportional to
$\prod_{p\in P(V_x)}e^{-\beta((d\theta)_p+2\pi k_p)^2/2}$, which does not
involve $k_{P\setminus P(V_x)}$; summing over $k_{P(V_x)}$ gives the density
$\prod_{p\in P(V_x)}\varphi_\beta((d\theta)_p)$ of the conditional law of
$\theta_{V_x}$ given $\theta_{E\setminus V_x}$ under $\mu_\beta$. Hence $\E[f\mid\theta_{E\setminus V_x},k_{P\setminus P(V_x)}]=
\mathcal P_xf$ for $f=f(\theta)$. Since $\sigma\bigl((e^{iF_p})_{p\notin
P(V_x)}\bigr)\subset\sigma(F_{P\setminus P(V_x)})\subset\sigma(\theta_{E
\setminus V_x},k_{P\setminus P(V_x)})$ and $\mathcal P_xf$ is measurable
with respect to the smallest of these, the tower property gives the two
remaining equalities.
\end{proof}

\subsubsection{The dual form of the cube estimate}\label{rM3:dualsec}

\begin{theorem}[dual form of $\mathrm{BP}(n,\varpi)$ on the plaquette
algebra]\label{rM3:dualform}
Let $P_\beta$ be the law of $F=d\theta+2\pi k$ under $\tilde\mu_\beta$
(Proposition~\ref{rM3:lawF}), and for $x\in X$ let $F_x^{\rm ext}:=
(F_p)_{p\notin P(V_x)}$. The inequality $\varpi\Var_\mu(f)\le\mathcal E_n
(f)$ holds for every real $f\in\mathcal A_F$ if and only if
\[
\varpi\Var_{P_\beta}(g)\le\frac1{n^3}\sum_{x\in X}\E_{P_\beta}\Bigl[\Var_{P_
\beta}\bigl(g\bigm|F_x^{\rm ext}\bigr)\Bigr]
\]
for every real $g\in L^2(P_\beta)$ with $g(F+2\pi j)=g(F)$ for all
$j\in\Z^P$.
\end{theorem}

\begin{proof}
The map $g\mapsto g(F)$ is an isometric isomorphism of the real
$2\pi$-periodic functions in $L^2(P_\beta)$ onto the real part of
$\mathcal A_F$ (a function of $(e^{iF_p})_p$ is a periodic function of
$F$ and conversely), with $\Var_\mu(g(F))=\Var_{P_\beta}(g)$ since
$\mu_\beta$ is the $\theta$-marginal of $\tilde\mu_\beta$. By
Corollary~\ref{rM3:exterior}, $\mathcal P_x(g(F))=\E[g(F)\mid F_x^{\rm
ext}]$, so $\mu((f-\mathcal P_xf)^2)=\E[(g-\E[g\mid F_x^{\rm ext}])^2]=
\E[\Var(g\mid F^{\rm ext}_x)]$.
\end{proof}

The right side is the Dirichlet form of a dynamics on the field strength
alone: resample the $3n^2(n+2)$ plaquette values inside the dual block
from their conditional law given the exterior values. The next result
describes that conditional law. Let $\Lambda_x:=\Z^{P(V_x)}\cap D_x$ and
$\mathcal M_x:=2\pi(1-\Pi_x)\Z^{P(V_x)}$; as in Lemma~\ref{rM3:lattices},
$\Lambda_x$ is a full-rank lattice in $D_x$ and $\mathcal M_x$ a full-rank
lattice in $D_x^\perp\cap\R^{P(V_x)}$, and $\Gamma_x=D_x\oplus\mathcal M_x$.

\begin{proposition}[the interior field given the exterior]
\label{rM3:interior}
Fix $x\in X$ and write $F^{\rm int}:=(F_p)_{p\in P(V_x)}$. Under
$\tilde\mu_\beta$, conditionally on $F^{\rm ext}_x$ (and also
conditionally on the finer data $(\theta_{E\setminus V_x},k_{P\setminus
P(V_x)})$: the two conditional laws coincide):
\begin{enumerate}[label=\textup{(\roman*)},leftmargin=*]
\item $F^{\rm int}$ lies in the coset $b+\Gamma_x$ of
Lemma~\ref{rM3:coset}, which is a function of $F^{\rm ext}_x$, and its
conditional law has density proportional to $e^{-\beta\lVert F^{\rm int}
\rVert^2/2}$ with respect to the Haar measure $\lambda_{D_x}\otimes
\mathrm{counting}$ of that coset;
\item $\Pi_xF^{\rm int}$ is Gaussian with mean $0$ and covariance
$\beta^{-1}\Pi_x$, whatever the exterior;
\item $(1-\Pi_x)F^{\rm int}$ lies in the coset $a_0+\mathcal M_x$, $a_0=
(1-\Pi_x)b$, a function of $F^{\rm ext}_x$, with conditional probabilities
proportional to $e^{-\beta\lVert a\rVert^2/2}$, $a\in a_0+\mathcal M_x$;
\item $\Pi_xF^{\rm int}$ and $(1-\Pi_x)F^{\rm int}$ are conditionally
independent.
\end{enumerate}
\end{proposition}

\begin{proof}
Given $(\theta_{E\setminus V_x},k_{P\setminus P(V_x)})$, the pair
$(\eta,K):=(\theta_{V_x},k_{P(V_x)})$ has density proportional to
$e^{-\beta\lVert b+d_x\eta+2\pi K\rVert^2/2}$ on $(-\pi,\pi]^{V_x}\times
\Z^{P(V_x)}$, with $b=b(\theta)$, and $F^{\rm int}=b+d_x\eta+2\pi K$. The
computation in the proof of Proposition~\ref{rM3:lawF}, with $E,P,d,
\mathcal B,\Lambda_{\mathcal B},\mathcal M$ replaced by $V_x,P(V_x),d_x,
D_x,\Lambda_x,\mathcal M_x$ and $\Phi$ replaced by $\Phi(b+\cdot)$, gives
for Borel $\Phi\ge0$
\begin{multline*}
\E\bigl[\Phi(F^{\rm int})\bigm|\theta_{E\setminus V_x},k_{P\setminus P(V_x)}
\bigr]\\
=\frac1{N(b)}\sum_{a\in\mathcal M_x}\int_{D_x}\Phi(b+g+a)\,
e^{-\beta\lVert b+g+a\rVert^2/2}\,\lambda_{D_x}(dg),
\end{multline*}
with $N(b)$ the normalisation. The measure
\[
\sum_{a\in\mathcal M_x}\int_{D_x}\delta_{b+g+a}\,\lambda_{D_x}(dg)
\]
is the Haar measure of the coset
$b+D_x+\mathcal M_x=b+\Gamma_x$, and the right side depends on $b$ only
through this coset. By Lemma~\ref{rM3:coset} the coset is a function of
$F^{\rm ext}_x$, so the displayed conditional law is also the conditional
law given $F^{\rm ext}_x$; this is (i). For (ii)--(iv) write
\[
b+g+a=(\Pi_xb+g)+((1-\Pi_x)b+a)
\]
with $\Pi_xb+g\in D_x$ and $(1-\Pi_x)b+a\in
D_x^\perp$, so
\[
\lVert b+g+a\rVert^2=\lVert\Pi_xb+g\rVert^2+\lVert(1-\Pi_x)b+a\rVert^2
\]
and the conditional law factorises into the image of the measure
\[
e^{-\beta\lVert\Pi_xb+g\rVert^2/2}\lambda_{D_x}(dg)
\]
under $g\mapsto\Pi_xb+g$, which is the centred Gaussian $N(0,\beta^{-1}\Pi_x)$ on $D_x$ by translation
invariance of $\lambda_{D_x}$, and the discrete law of (iii). The coset
$a_0+\mathcal M_x=(1-\Pi_x)(b+\Gamma_x)$ is a function of $F^{\rm ext}_x$.
\end{proof}

\begin{corollary}[the cube photon is never screened]\label{rM3:photon}
For $c\in\R^P$ and every $x\in X$,
\[
\E\bigl[\Var(\ip cF\mid F^{\rm ext}_x)\bigr]\ \ge\ \beta^{-1}\lVert\Pi_x
c^{\rm int}\rVert^2,\qquad c^{\rm int}:=c|_{P(V_x)},
\]
with equality if and only if $\ip{c^{\rm int}}{(1-\Pi_x)F^{\rm int}}$ is
a.s. a function of $F^{\rm ext}_x$. The lower bound is the Gaussian cube
response of Lemma~\ref{rD:gaussresp} in three dimensions.
\end{corollary}

\begin{proof}
Decompose
\[
\ip cF=\ip{c^{\rm ext}}{F^{\rm ext}_x}+\ip{c^{\rm int}}{\Pi_xF^{\rm int}}
+\ip{c^{\rm int}}{(1-\Pi_x)F^{\rm int}};
\]
the first term is
$F^{\rm ext}_x$-measurable, the other two are conditionally independent by
Proposition~\ref{rM3:interior}(iv), and the conditional variance of the
second is $\beta^{-1}\lVert\Pi_xc^{\rm int}\rVert^2$ by (ii).
\end{proof}

\begin{proposition}[the dual action of $\mathcal P_x$ on characters]
\label{rM3:characters}
For $m\in\Z^P$ let $\chi_m(F)=e^{i\ip mF}$, so that the $\chi_m$ span a dense
subspace of $\mathcal A_F$. Write $m=m^{\rm int}+m^{\rm ext}$ according to
$P=P(V_x)\sqcup(P\setminus P(V_x))$. Then
\begin{gather*}
\mathcal P_x\chi_m=\chi_{m^{\rm ext}}\,\frac{\Sigma(m^{\rm int})}
{\Sigma(0)},\\
\Sigma(r)=\sum_{\substack{m'\in\Z^{P(V_x)}\\ \delta(r+m')|_{V_x}=0}}
e^{-\lVert m'\rVert^2/2\beta}\,e^{i\ip{\delta(r+m')}{\theta}_{E\setminus V_x}},
\end{gather*}
where $\ip\cdot\cdot_{E\setminus V_x}$ is the sum over exterior links; by
Lemma~\ref{rM3:coset} each exponential $e^{i\ip{\delta\tilde m}\theta_{E
\setminus V_x}}$ with $\delta\tilde m|_{V_x}=0$ is a function of $F^{\rm
ext}_x$.
\end{proposition}

\begin{proof}
Since $m^{\rm int}$ is integer, $\chi_{m^{\rm int}}(F^{\rm int})=e^{i\ip
{m^{\rm int}}{b+d_x\eta}}$ does not depend on $K$, and summing the
conditional density over $K$ gives $\prod_{p\in P(V_x)}\varphi_\beta(b_p+(d_x
\eta)_p)$. Insert the Poisson summation form $\varphi_\beta(t)=(2\pi
\beta)^{-1/2}\sum_{m'\in\Z}e^{-m'^2/2\beta}e^{im't}$ of the Villain
weight (the Fourier coefficients of the $2\pi$-periodisation of $e^{-\beta
t^2/2}$), exchange the absolutely convergent sum with the integral, and
integrate over $\eta\in(-\pi,\pi]^{V_x}$: $\int e^{i\ip{\tilde m}{d_x\eta}}
d\eta/(2\pi)^{\lvert V_x\rvert}=\mathbf 1\{\delta\tilde m|_{V_x}=0\}$, and for
such $\tilde m$, $\ip{\tilde m}b=\ip{\tilde m}{d\theta^0}=\ip{\delta\tilde m}
{\theta^0}=\ip{\delta\tilde m}\theta_{E\setminus V_x}$. The numerator is this
with $\tilde m=m^{\rm int}+m'$, the denominator with $\tilde m=m'$.
\end{proof}

\begin{remark}[the dual height model and the sine-Gordon field]
\label{rM3:nodual}
Poisson resummation of every $\varphi_\beta$ and integration over $\theta$
give
\[
Z=\Bigl(\frac1{2\pi\beta}\Bigr)^{\lvert P\rvert/2}\sum_{m\in\Z^P,\ \delta m
=0}e^{-\lVert m\rVert^2/2\beta},
\]
a sum over integer-valued coclosed $2$-cochains. They include the $m=
\delta h$, $h\in\Z^C$, the gradients of integer-valued height functions
on the dual lattice, and the restriction of the sum to these is the
partition function of an integer-valued height model; duality
transformations of this kind are the starting point of~\cite{GM82}
(and, in four dimensions, of~\cite{FS82}), see~\cite{RFN24}. Proposition~\ref{rM3:characters} expresses the cube
move on characters as a ratio of two sums of this type over the
$m'$ supported in the dual block, weighted by the phases
$e^{i\ip{\delta\tilde m}\theta}$, which are functions of the exterior field
strength. These weights are complex in general, so the formula is not
the conditional expectation of a block of the height model, and we do not
know a positive dual measure whose block resampling is intertwined with the
cube move. Likewise the Gaussian (sine-Gordon type) field that arises
from a further Gaussian transformation of the Coulomb energy of
Remark~\ref{rM3:coulomb}, as in the approach of~\cite{GM82} described
in~\cite{RFN24}, enters as a variable in a representation of
expectations; it is not constructed here as a random variable on the
probability space of $\tilde\mu_\beta$ coupled to $F$. So
Theorem~\ref{thm:ffscalar} applied to a massive scalar of mass $m_D$ gives
no statement about $\mu_\beta$, and in this subsection the monopole
screening enters $\mathrm{BP}(n,\varpi)$ on $\mathcal A_F$ only through
the inequality of Theorem~\ref{rM3:dualform}. Theorem~\ref{rM3:DH} below makes this
precise in the Gaussian approximation, where the cube move does act on
the field strength itself.
\end{remark}
\subsubsection{Single plaquettes: the photon gives a $\beta$-independent
constant}\label{rM3:plaqsec}

\begin{lemma}[second moment of the field strength]\label{rM3:moment}
Under $\tilde\mu_\beta$, $\E\lVert M\rVert^2\le\dim\mathcal B^\perp/\beta
=(L^3+2)/\beta$, $\E\lVert G\rVert^2=\dim\mathcal B/\beta=(2L^3-2)/\beta$,
and $\E F_p^2\le1/\beta$ for every plaquette $p$.
\end{lemma}

\begin{proof}
$\ker d\subset\R^E$ consists of the exact $1$-cochains ($\dim L^3-1$) and
the harmonic ones ($\dim3$), so $\dim\mathcal B=3L^3-(L^3+2)=2L^3-2$ and
$\dim\mathcal B^\perp=L^3+2=:N$. By Proposition~\ref{rM3:lawF}(ii), $G$ is
centred Gaussian with covariance $\beta^{-1}\Pi$, so $\E\lVert G\rVert^2=
\beta^{-1}\operatorname{tr}\Pi=\dim\mathcal B/\beta$. Let $\mathcal M^*=\{w
\in\mathcal B^\perp:\ip wa\in\Z\ \forall a\in\mathcal M\}$ be the dual lattice
and $\lvert\mathcal M\rvert$ the covolume (Lemma~\ref{rM3:lattices}(iii)).
The Poisson summation formula for the Gaussian $e^{-\beta\lVert y\rVert^2/2}$
on $\mathcal B^\perp\cong\R^N$ gives, for $\beta>0$,
\begin{gather*}
Z_{\mathcal M}(\beta):=\sum_{a\in\mathcal M}e^{-\beta\lVert a\rVert^2/2}=
\frac{(2\pi/\beta)^{N/2}}{\lvert\mathcal M\rvert}\,\Theta(\beta),\\
\Theta(\beta)=\sum_{w\in\mathcal M^*}e^{-2\pi^2\lVert w\rVert^2/\beta}.
\end{gather*}
Both series and their $\beta$-derivatives converge uniformly on compact
subsets of $(0,\infty)$, and every term of $\Theta$ is nondecreasing in
$\beta$, so $\Theta'\ge0$. Hence
\[
\E\lVert M\rVert^2=-2\frac{d}{d\beta}\log Z_{\mathcal M}(\beta)=\frac N\beta
-2\frac{\Theta'(\beta)}{\Theta(\beta)}\le\frac N\beta .
\]
Since $G$ and $M$ are independent and $\E G=0$, $\E\lVert F\rVert^2=
\E\lVert G\rVert^2+\E\lVert M\rVert^2\le3L^3/\beta$. The translations of
$\Lambda_L$ and the cyclic permutation of the coordinate axes act on
$(\theta,k)$ preserving $\tilde\mu_\beta$, and they act on $\R^P$ by
permutations of the plaquettes up to signs; this group acts transitively on
$P$, so $\E F_p^2$ does not depend on $p$ and equals $\E\lVert F\rVert^2/
(3L^3)\le1/\beta$.
\end{proof}

\begin{proposition}[odd plaquette functions]\label{rM3:plaq}
Let $\beta>1$, $n\ge2$, $L\ge n+2$, and let $p$ be a plaquette. Then
$f=\sin F_p=\operatorname{Im}e^{i(d\theta)_p}\in\mathcal A_F$ satisfies
\[
\mathcal E_n(f)\ge\frac{(1+2/n)(1-1/\beta)}4\,\Var_\mu(f),
\]
and the constant is at least $\frac18$ for $\beta\ge2$. The same holds for
$c\sin F_p$, $c\in\R$.
\end{proposition}

\begin{proof}
$\Var_\mu(f)\le\E\sin^2F_p\le\E F_p^2\le1/\beta$ by Lemma~\ref{rM3:moment}.
If $p\notin P(V_x)$, $f$ is $F^{\rm ext}_x$-measurable and contributes
nothing; $p\in P(V_x)$ for exactly $n\cdot n(n+2)=n^2(n+2)$ points $x$ (the
number of plaquettes of $P(V_0)$ in the plane of $p$,
Sec.~\ref{rM3:block}). Fix such an $x$. Then $p$ has a link $e\in V_x$; the
four plaquettes containing $e$ lie in $P(V_x)$, so $w=d_x\mathbf 1_e\in D_x$
has $\lVert w\rVert^2=4$ and $\ip{\mathbf 1_p}w=\pm1$, whence
\[
\sigma^2:=\beta^{-1}\lVert\Pi_x\mathbf 1_p\rVert^2\in\bigl[\tfrac1{4\beta},
\tfrac1\beta\bigr].
\]
Write $F_p=X+a$ with $X=\ip{\mathbf 1_p}{\Pi_xF^{\rm int}}$ and $a=\ip{\mathbf
1_p}{(1-\Pi_x)F^{\rm int}}$. By Proposition~\ref{rM3:interior}(ii),(iv),
conditionally on $F^{\rm ext}_x$ the variable $X$ is $N(0,\sigma^2)$ and
independent of $a$. For $X\sim N(0,\sigma^2)$ and a constant $a$,
$\E\sin(X+a)=e^{-\sigma^2/2}\sin a$ and $\E\sin^2(X+a)=\frac12(1-e^{-2
\sigma^2}\cos2a)$, so
\begin{align*}
V(a):=\Var\sin(X+a)&=\frac{1-e^{-2\sigma^2}}2-\bigl(e^{-\sigma^2}-e^{-2\sigma^2}
\bigr)\sin^2a\\
&\ge\sigma^2-\sigma^4-\sigma^2a^2,
\end{align*}
using $1-e^{-u}\ge u-u^2/2$ and $e^{-\sigma^2}(1-e^{-\sigma^2})\le\sigma^2$.
By the law of total variance, $\Var(f\mid F^{\rm ext}_x)\ge\E[V(a)\mid
F^{\rm ext}_x]$ (given $(a,F^{\rm ext}_x)$ the variable $X$ is still
$N(0,\sigma^2)$, so $\Var(f\mid a,F^{\rm ext}_x)=V(a)$). Moreover $\E X^2=
\sigma^2$ and $\E[Xa]=0$ (condition on $(a,F^{\rm ext}_x)$), so $\E a^2=
\E F_p^2-\sigma^2\le1/\beta-\sigma^2$. Hence
\begin{align*}
\E\bigl[\Var(f\mid F^{\rm ext}_x)\bigr]&\ge\sigma^2-\sigma^4-\sigma^2
\Bigl(\frac1\beta-\sigma^2\Bigr)\\
&=\sigma^2\Bigl(1-\frac1\beta\Bigr)\ge\frac{1-1/\beta}{4\beta}.
\end{align*}
By Corollary~\ref{rM3:exterior}, $\mu((f-\mathcal P_xf)^2)=\E[\Var(f\mid
F^{\rm ext}_x)]$, and summing over the $n^2(n+2)$ points $x$,
$\mathcal E_n(f)\ge n^{-3}n^2(n+2)(1-1/\beta)/(4\beta)\ge\frac{(1+2/n)(1-
1/\beta)}4\Var_\mu(f)$. For $\beta\ge2$ the constant is at least
$(1+2/n)/8\ge\frac18$. Both sides are quadratic in $c$.
\end{proof}

\begin{remark}[what this shows, and the even functions]\label{rM3:plaqremark}
For $\SU(2)$, Proposition~\ref{rD:plaquette} obtains the inequality of
$\mathrm{BP}$ for plaquette traces with a constant of order
$1/(\beta\log\beta)$, and the second-moment bound of
Theorem~\ref{rDtwo:secondmoment} removes the logarithmic loss
(Corollary~\ref{rDtwo:bpplaquette}, Remark~\ref{rD:plaqloss}). For the odd
functions of one Villain plaquette the constant is independent of $\beta$
and $L$ and needs only a second moment: the
exact part of the interior field is a Gaussian of variance $\asymp\beta^{-1}$
whatever the exterior (Proposition~\ref{rM3:interior}), and the second
moment of $F_p$ is at most $\beta^{-1}$ (Lemma~\ref{rM3:moment}). For the
even function $\cos F_p$ the same argument gives $\E\Var(\cos F_p\mid
F^{\rm ext}_x)\ge\E[(1-e^{-\sigma^2})^2/2]\ge c\beta^{-2}$ for $\beta\ge1$
(the conditional variance of $\cos(X+a)$ is $\frac12(1-e^{-\sigma^2})^2+
(e^{-\sigma^2}-e^{-2\sigma^2})\sin^2a$), and a $\beta$-independent constant
would follow from $\Var(\cos F_p)\le C\beta^{-2}$. For this
$\E(1-\cos F_p)^2\le\frac14\E F_p^4\le C\beta^{-2}$ would suffice, and
since $F_p=G_p+M_p$ with $G_p$ Gaussian of variance at most $\beta^{-1}$, a
fourth-moment bound $\E M_p^4\le C'\beta^{-2}$ on the monopole component
would give it; Lemma~\ref{rM3:moment} gives only the second moment, and
this fourth-moment bound is not proved here. By the
symmetry $(\theta,k)\mapsto(-\theta,-k)$, which preserves $\tilde\mu_\beta$
and the exterior $\sigma$-algebras, $\cos F_p$ and $\sin F_p$ are
orthogonal both for the variance and for the Dirichlet form, so the
inequality on $\operatorname{span}\{\cos F_p,\sin F_p\}$ is equivalent to
the two separate inequalities. These are statements at the scale of one
plaquette. Their mechanism is the photon; the monopoles enter only
through the second moment of Lemma~\ref{rM3:moment}. They say
nothing about the long-wavelength functions that decide
$\mathrm{BP}(n,\varpi)$ (Theorem~\ref{rD:maxwell}).
\end{remark}

\subsubsection{The layer flux: where the monopoles break the Maxwell
mechanism}\label{rM3:layersec}

Fix $x$ and the plane $\{1,2\}$. For $t\in x_3+[0,n)$ let $Q_t$ be the
$(n+1)^2$ plaquettes $(y,12)$ with $y_3=t$ and $(y_1,y_2)\in(x_1,x_2)+[-1,
n-1]^2$, $q_t\in Q_t$ the one with $(y_1,y_2)=(x_1-1,x_2-1)$, $R_t=Q_t
\setminus\{q_t\}$, $s_t=x_3+n-t\in[1,n]$, $Q'_t=Q_t+s_te_3$, and let $T_t$ be
the set of the $(n+1)^2s_t$ cubes $c_{z+je_3}$, $(z,12)\in Q_t$, $0\le j<s_t$.
By Sec.~\ref{rM3:block}, the plaquettes of $P(V_x)$ in the plane $\{1,2\}$
are the disjoint union of the $R_t$ (this uses $L\ge n+2$).

\begin{lemma}[layer flux]\label{rM3:layer}
Let $L\ge n+2$. There are integer $2$-chains $W_t$ supported on plaquettes
not in $P(V_x)$ such that $Q'_t\cap P(V_x)=\emptyset$, $q_t\notin P(V_x)$,
and for every $F\in\R^P$
\begin{gather*}
\Phi_t(F):=\sum_{p\in R_t}F_p=E_t(F)-\ip{dF}{\mathbf 1_{T_t}},\\
E_t(F):=\sum_{p\in Q'_t}F_p+\ip F{W_t}-F_{q_t}.
\end{gather*}
For the Villain field strength $F=d\theta+2\pi k$, $\ip{dF}{\mathbf 1_{T_t}}
=2\pi N_t$ with $N_t=\sum_{c\in T_t}\mathsf n_c\in\Z$ the monopole number
of $T_t$, and $E_t(F)$ is a function of $F^{\rm ext}_x$.
\end{lemma}

\begin{proof}
By the formula for $d$ on $2$-cochains (Sec.~\ref{rM3:setting}), the
boundary chain $\partial c_z:=d^{\mathsf T}\mathbf 1_{c_z}$ is $\sum_\rho
\sigma_\rho(\mathbf 1_{(z+e_\rho,\mu\nu)}-\mathbf 1_{(z,\mu\nu)})$, $\{\mu,
\nu\}$ the complement of $\rho$, $\sigma_1=\sigma_3=1$, $\sigma_2=-1$. In
$\partial\mathbf 1_{T_t}$ the $12$-faces telescope in $j$ to $\mathbf 1_{Q'_t}-
\mathbf 1_{Q_t}$. For fixed $j$, the $23$-faces $\pm[(z+e_1,23)-(z,23)]$
telescope in $z_1\in x_1+[-1,n-1]$ and leave faces with first coordinate
$x_1-1$ or $x_1+n$; a $23$-plaquette has its three base points $y,y+e_2,
y+e_3$ with the same first coordinate, which is then not in $x_1+[0,n)$, so
these faces are not in $P(V_x)$. The $13$-faces are treated in the same way
with the second coordinate. Put $W_t=\partial\mathbf 1_{T_t}-(\mathbf 1_{Q'_t}
-\mathbf 1_{Q_t})$. The plaquettes of $Q'_t$ have all base points with third
coordinate $x_3+n$, and $q_t\notin P(V_x)$ by the description of
$P(V_x)$ in Sec.~\ref{rM3:block}. Since $L\ge n+2$, the residues
$x_i-1,\dots,x_i+n$ are distinct, so the cubes of $T_t$ are distinct and
the cancellations are as stated. Now $\ip F{\partial\mathbf 1_{T_t}}=\ip{dF}
{\mathbf 1_{T_t}}$ gives $\sum_{Q'_t}F-\sum_{Q_t}F+\ip F{W_t}=\ip{dF}{\mathbf
1_{T_t}}$, which is the identity after subtracting $F_{q_t}$. For the
Villain field, $dF=2\pi dk=2\pi\mathsf n$, and $E_t$ involves only
plaquettes outside $P(V_x)$.
\end{proof}

For the Maxwell field $dF=0$ and $\Phi_t$ is exterior-measurable, which is
the whole mechanism of Theorem~\ref{rD:maxwell}: a cube move cannot change
the flux through the layers it reads. For the Villain field the move
changes $\Phi_t$ exactly by $2\pi$ times the change of the monopole number
of $T_t$, a region of the dual block. The layer flux is relaxed by the
cube dynamics only through monopoles created or annihilated inside the
block, and with weight $2\pi$, an integer multiple of the period. A
function of the plaquette angles $e^{iF_p}$ sees $\Phi_t$ only modulo
$2\pi$ through products of characters, and $e^{i\Phi_t}=e^{iE_t}$ is
exterior-measurable. So the long-wavelength test functions of
Theorem~\ref{rD:maxwell} must be built from non-periodic functions of
$F$, such as $\sin F_p$ summed over many plaquettes, whose layer sums are
not of the form $e^{i\Phi_t}$. Theorem~\ref{rM3:DH} shows what happens when
the integer constraint on $dF$ is replaced by a Gaussian penalty.
\subsubsection{The Debye--H\"uckel Gaussian field}\label{rM3:DHsec}

By Proposition~\ref{rM3:lawF}, $P_\beta$ has density $\propto e^{-\beta
\lVert F\rVert^2/2}$ with respect to $\lambda_{\mathcal B}\otimes
\mathrm{counting}$ on the union $\mathcal B+\mathcal M$ of translates of
$\mathcal B$, and every $F$ in its support has $dF\in2\pi\Z^C$. In the dilute-gas (Debye--H\"uckel) approximation the
integer charge $\mathsf n=dF/2\pi$ is replaced by a centred Gaussian with
variance $2z$ per cube, that is, the constraint by the weight
$e^{-\lVert dF\rVert^2/(16\pi^2z)}=e^{-\beta\lVert dF\rVert^2/(2\mathsf m^2)}$
with $\mathsf m^2=8\pi^2\beta z$ (Remark~\ref{rM3:coulomb}). We take $\mathsf
m>0$ as a free parameter.

\begin{definition}[the Debye--H\"uckel field]\label{rM3:DHdef}
For $\beta,\mathsf m>0$ let $Q_{\mathsf m}=\beta(I+\mathsf m^{-2}\delta d)$
on $\R^P$, a positive definite operator ($Q_{\mathsf m}\ge\beta$), and let
$\gamma^{\rm DH}$ be the centred Gaussian law on $\R^P$ with covariance
$Q_{\mathsf m}^{-1}$. For $x\in X$ let $\mathcal P^{\rm DH}_x$ be the
conditional expectation given $(F_p)_{p\notin P(V_x)}$, and
$\varpi_{\rm DH}(n,\mathsf m)$ the largest $\varpi$ with
\[
\varpi\Var_{\gamma^{\rm DH}}(f)\le\frac1{n^3}\sum_{x\in X}\bigl\lVert f-
\mathcal P^{\rm DH}_xf\bigr\rVert^2_{L^2(\gamma^{\rm DH})}
\]
for all real $f\in L^2(\gamma^{\rm DH})$,
which is $\varpi(Q_{\mathsf m};(P(V_x))_x)$ of~\eqref{eq:ffvarpi} with the
normalisation $n^{-3}$.
\end{definition}

On the exact cochains $\mathcal B$, $Q_{\mathsf m}=\beta$, and on $\delta
\R^C$, $Q_{\mathsf m}\delta u=\beta\delta(1+\mathsf m^{-2}(-\Delta_3))u$, so
the covariance is $Q_{\mathsf m}^{-1}=\beta^{-1}[\Pi+H+\mathsf m^2(\mathsf m^2
+\delta d)^{-1}\Pi_c]$, $\Pi_c$ and $H$ the projections onto $\delta\R^C$
and $\mathcal H$. The
exact part is the same massless Maxwell field as in
Proposition~\ref{rM3:lawF}; nevertheless the off-diagonal covariance
$\ip{\mathbf 1_p}{Q_{\mathsf m}^{-1}\mathbf 1_q}=-\beta^{-1}\ip{d\mathbf 1_p}
{(\mathsf m^2-\Delta_3)^{-1}d\mathbf 1_q}$, $p\ne q$ (from $(I+\mathsf m^{-2}
\delta d)^{-1}=I-\delta(\mathsf m^2+d\delta)^{-1}d$ and $d\delta=-\Delta_3$
on $\R^C$), decays exponentially at rate $\asymp\mathsf m$ for small
$\mathsf m$, as the lattice resolvent $(\mathsf m^2-\Delta_3)^{-1}$ does.
The field strength of
$\gamma^{\rm DH}$ is massive although its exact part is not, and the
mechanism is the one of the Villain decomposition: the coexact part cancels
the long-range tail of $\Pi$.

\begin{theorem}[the cube condition for the Debye--H\"uckel field]
\label{rM3:DH}
Let $n\ge2$, $L\ge n+2$, $\beta,\mathsf m>0$, $s=n\mathsf m$, and
$\alpha_n,\rho_n$ as in Lemma~\ref{lem:ffparabola}. Then
\begin{enumerate}[label=\textup{(\roman*)},leftmargin=*]
\item $\displaystyle\varpi_{\rm DH}(n,\mathsf m)\ge\frac1{\alpha_n^3}\,
\frac{s^2}{2s^2+60\rho_n}\ge\Bigl(\frac56\Bigr)^3\frac{s^2}{2s^2+60}$;
\item $\displaystyle\varpi_{\rm DH}(n,\mathsf m)\le\min\Bigl\{1+\frac2n,\
\frac{b_n^2s^2}2\Bigr\}$, where
\[
b_n^2=\frac{(n+1)^3(2n+1)}{3n^4}\le\frac{45}{16};
\]
\item the slowly varying wave $f=\sum_y\cos(2\pi y_1/L)F_{(y,12)}$ satisfies
$\displaystyle\frac{\mathcal E^{\rm DH}_n(f)}{\Var(f)}\le\Bigl(\frac{2\pi
\sqrt{2n(n+2)}}L+b_ns\Bigr)^2$, $\mathcal E^{\rm DH}_n$ the right side
of the inequality in Definition~\ref{rM3:DHdef}.
\end{enumerate}
The bounds are independent of $\beta$, and (i), (ii) are independent of
$L$. In particular $c\min(1,s^2)\le\varpi_{\rm DH}\le\min(1+\frac2n,
\frac{45}{32}s^2)$ with $c=(5/6)^3/62$.
\end{theorem}

\begin{proof}
(i) For a plaquette $p=(y,\mu\nu)$ and $x\in X$ put $\chi_x(p)=\prod_{i=1}^3
\tau(y_i-x_i)$, $\tau$ the parabolic profile. If $\chi_x(p)\ne0$ then $y\in x
+[0,n)^3$, so the link $(y,\mu)$ of $p$ lies in $V_x$ and $p\in P(V_x)$; and
$\sum_x\chi_x(p)=1$, $\sum_x\chi_x(p)^2=S_2^3$ as in~\eqref{eq:ffchisums}.
For $v\in\R^P$ put $v_x=\chi_xv$, supported in $P(V_x)$, with $\sum_xv_x=v$.
For $\rho\in\{1,2,3\}$ let $w_\rho(y)=\sigma_\rho v_{(y,\mu\nu)}$, $\{\mu,\nu\}$
the complement of $\rho$; then $\lVert v\rVert^2=\sum_\rho\lVert w_\rho
\rVert^2$, $(dv)_z=\sum_\rho\nabla_\rho w_\rho(z)$ by the formula for $d$,
and, since $\chi_x(p)$ depends only on the base point of $p$,
\[
(d(\chi_xv))_z=\chi_x(z)(dv)_z+\sum_\rho(\nabla_\rho\chi_x)(z)\,w_\rho(z+e_\rho),
\]
where $\chi_x(z)=\prod_i\tau(z_i-x_i)$. By $(a+b_1+b_2+b_3)^2\le2a^2+6\sum
b_\rho^2$ and $\sum_x\chi_x(z)^2=S_2^3$, $\sum_x(\nabla_\rho\chi_x(z))^2=
S_2^2D_2$ (the sum factorises over the coordinates, and $L\ge n+1$),
\[
\sum_x\lVert d(\chi_xv)\rVert^2\le2S_2^3\lVert dv\rVert^2+6S_2^2D_2\lVert v
\rVert^2 .
\]
Hence
\begin{align*}
\sum_x\ip{v_x}{Q_{\mathsf m}v_x}&=\beta\sum_x\Bigl(\lVert\chi_xv\rVert^2+\frac{
\lVert d(\chi_xv)\rVert^2}{\mathsf m^2}\Bigr)\\
&\le C_0\,\beta\Bigl(\lVert v
\rVert^2+\frac{\lVert dv\rVert^2}{\mathsf m^2}\Bigr)=C_0\ip v{Q_{\mathsf m}v}
\end{align*}
with $C_0=S_2^2\max(S_2+6D_2\mathsf m^{-2},2S_2)\le S_2^2(2S_2+6D_2\mathsf
m^{-2})$. Lemma~\ref{lem:fflions}, which is stated for an arbitrary finite
family of coordinate sets, applied with $T=\R^P$, the sets $P(V_x)$ and
$d=3$, gives $\varpi_{\rm DH}\ge1/(n^3C_0)$. With $S_2=\alpha_n/n$, $D_2=
\beta_n/n^3$ and $\beta_n=10\rho_n\alpha_n$ (here $\beta_n$ is the constant
of Lemma~\ref{lem:ffparabola}, not the coupling),
\[
n^3C_0\le2\alpha_n^3+\frac{6\alpha_n^2\beta_n}{s^2}=\alpha_n^3\Bigl(2+\frac{60
\rho_n}{s^2}\Bigr),
\]
which is the first bound; the second uses $\alpha_n\le6/5$, $\rho_n<1$.

(ii), (iii) The bound $1+2/n$: for $f=F_p$, $\lVert f-\mathcal P^{\rm DH}_x
f\rVert^2\le\Var(F_p)$, and it vanishes unless $p\in P(V_x)$, which holds
for $n^2(n+2)$ points $x$. For the other two bounds let $c\in\R^P$ vanish
off the plaquettes $(y,12)$ and take on them either the values $c_{(y,12)}=
1$ (case (ii)) or $c_{(y,12)}=\cos(2\pi y_1/L)$ (case (iii)), and put $f=
\ip cF$. In both cases $dc=0$: by the formula for $d$ on $2$-cochains only
the term $\rho=3$, $c_{(z+e_3,12)}-c_{(z,12)}$, can contribute, and $c_{(y,
12)}$ does not depend on $y_3$. Hence $Q_{\mathsf m}c=\beta c$ and $\Var(f)=
\ip c{Q_{\mathsf m}^{-1}c}=\lVert c\rVert^2/\beta$, which is $L^3/\beta$ in
case (ii) and $L^3/(2\beta)$ in case (iii) (as $L\ge4$). Fix $x$, put $\bar
c=1$ in case (ii) and $\bar c=\cos(2\pi(x_1-1)/L)$ in case (iii),
$c^{\rm in}=c\mathbf 1_{P(V_x)}$,
$c^{\rm ext}=c-c^{\rm in}$, and with Lemma~\ref{rM3:layer}
\[
g=\ip{c^{\rm ext}}F+\bar c\sum_tE_t(F),
\]
a function of $(F_p)_{p\notin P(V_x)}$. Since $c^{\rm in}$ is supported on
$\bigsqcup_tR_t$ and $\sum_tE_t=\sum_t\Phi_t+\ip{dF}\vartheta$ with $\vartheta=\sum_t
\mathbf 1_{T_t}$,
\[
f-g=\ip hF-\bar c\ip F{\delta\vartheta},\qquad h=c^{\rm in}-\bar c\sum_t\mathbf 1_{R_t}.
\]
In case (ii), $h=0$. In case (iii), as in Theorem~\ref{rD:maxwell}, $h$
is supported on the $n^2(n+2)$ plaquettes of the layers, and on $(y,12)\in
R_t$, $y_1=x_1-1+j$, $0\le j\le n$, $\lvert h_{(y,12)}\rvert\le2\pi j/L\le2\pi
n/L$, so $\E\ip hF^2=\ip h{Q_{\mathsf m}^{-1}h}\le\lVert h\rVert^2/\beta\le
4\pi^2n^4(n+2)/(\beta L^2)$, using $Q_{\mathsf m}\ge\beta$. Next, $u=(1+\mathsf m^{-2}(-\Delta_3))^{-1}\vartheta$ satisfies
$Q_{\mathsf m}\beta^{-1}\delta u=\delta\vartheta$, so
\begin{align*}
\E\ip F{\delta\vartheta}^2&=\ip{\delta\vartheta}{Q_{\mathsf m}^{-1}\delta\vartheta}=\frac1
\beta\Bigl\langle\vartheta,\frac{-\Delta_3}{1+\mathsf m^{-2}(-\Delta_3)}\vartheta
\Bigr\rangle\\
&\le\frac{\mathsf m^2}\beta\lVert\vartheta\rVert^2,
\end{align*}
since $\lambda/(1+\lambda/\mathsf m^2)\le\mathsf m^2$ for $\lambda\ge0$. The
cube $c_z$ lies in $T_t$ if and only if $(z_1,z_2)\in(x_1,x_2)+[-1,n-1]^2$
and $t\le z_3\le x_3+n-1$, so $\vartheta(z)=z_3-x_3+1$ on those cubes with $z_3\in
x_3+[0,n)$, and $\lVert\vartheta\rVert^2=(n+1)^2\sum_{j=1}^nj^2=(n+1)^3n(2n+1)/6$.
Since $g$ is measurable for the exterior, $I-\mathcal P^{\rm DH}_x$ is an
orthogonal projection, and $\lvert\bar c\rvert\le1$,
\[
\lVert f-\mathcal P^{\rm DH}_xf\rVert\le\lVert f-g\rVert\le\beta^{-1/2}
\bigl(a_h+\mathsf m\lVert\vartheta\rVert\bigr),
\]
with $a_h=0$ in case (ii) and $a_h=2\pi n^2\sqrt{n+2}/L$ in case (iii).
Summing the squares over the $L^3$ points $x$ and dividing by $n^3$: in
case (ii), $\varpi_{\rm DH}L^3/\beta\le n^{-3}L^3\mathsf m^2\lVert\vartheta
\rVert^2/\beta$, that is $\varpi_{\rm DH}\le\mathsf m^2(n+1)^3(2n+1)/(6n^2)=
b_n^2s^2/2$; in case (iii), $\mathcal E^{\rm DH}_n(f)/\Var(f)\le2n^{-3}(a_h+
\mathsf m\lVert\vartheta\rVert)^2=(2\pi\sqrt{2n(n+2)}/L+b_ns)^2$. Finally $b_n^2=
\frac13(1+\frac1n)^3(2+\frac1n)$ decreases in $n$ and equals $\frac{45}{16}$
at $n=2$.
\end{proof}

\begin{remark}[reading of Theorem~\ref{rM3:DH}]\label{rM3:DHremark}
With $n=\ell/a$ and $\mathsf m=m_{\rm phys}a$, $s=m_{\rm phys}\ell$, and the
theorem says that the Debye--H\"uckel field satisfies $\mathrm{BP}(n,
\varpi)$ with $\varpi\ge(5/6)^3s^2/(2s^2+60)$, uniformly in $L$, $\beta$ and
the lattice spacing, and with no better constant than
$\frac12b_n^2s^2\le\frac{45}{32}s^2$, on every torus. This is the behaviour of the free massive
scalar (Theorem~\ref{thm:ffscalar}, $\varpi\asymp s^2/(s^2+10d)$) and of the
Proca field (Theorem~\ref{thm:ffproca}, Proposition~\ref{prop:ffupper}),
with one difference: the mass term acts only on the coexact part, through
$\mathsf m^{-2}\lVert dF\rVert^2$, while the exact part keeps its massless
law. The proof of (i) shows how this suffices: cutting $v$ by the profile
creates the divergence $(\nabla\chi_x)\cdot w$, and that is paid by the
monopole term at the price $D_2/\mathsf m^2$, exactly as the mass term pays
the cut of a scalar field in Lemma~\ref{lem:ffpartition}. In the proofs of
(ii) and (iii) the test functions are no longer frozen, because the layer
flux $\Phi_t$ now differs from an exterior function by the Gaussian charge
$\ip{dF}{\mathbf 1_{T_t}}$, of variance $\le\mathsf m^2\lVert\mathbf
1_{T_t}\rVert^2/\beta$. The test function of (ii) is the flux through the
planes $y_3=\mathrm{const}$ of the torus, a harmonic direction that is
exactly frozen when $\mathsf m\to0$; that of (iii) is the slow wave of
Theorem~\ref{rD:maxwell}, and setting $\mathsf m=0$ in (iii) gives the
three-dimensional form of Theorem~\ref{rD:maxwell}. So in the Gaussian
approximation
the dynamically generated mass does produce the cube constant that
Remark~\ref{rem:rSmissing} asks for, on cubes of the size of the inverse
mass, and with a constant of order $(m_{\rm phys}\ell)^2$ on smaller cubes.
\end{remark}
\subsubsection{A necessary condition, and comparison}\label{rM3:compare}

\begin{proposition}[correlation length]\label{rM3:xi}
Let $\mathcal L$ be an infinite set of even $L\ge2n+4$ such that
$\mathrm{BP}(n,\varpi)$ holds for the Villain measure $\mu_\beta$ on
$\Lambda_L=(\Z/L\Z)^3$ for every $L\in\mathcal L$, and let $\omega$ be a
torus limit along $\mathcal L$ (defined as after
Theorem~\ref{thm:avgdecay}, with $\Z^3$ in place of $\Z^4$). If $f,g$ are bounded local gauge-invariant
functions with $\lvert\omega(f\tau_dg)-\omega(f)\omega(g)\rvert\ge c_0e^{-d/
\xi}$ for some $c_0>0$, $\xi>0$ and infinitely many $d$, then $\varpi\le e^2
q_nn/\xi$, $q_n=(2+1/n)^3$. With $na=\ell$ and $\xi=1/(m_{\rm phys}a)$
this reads $\varpi\le e^2q_n\,m_{\rm phys}\ell$.
\end{proposition}

\begin{proof}
The proof of Corollary~\ref{rD:xi} uses only the second bound of
Theorem~\ref{thm:avgdecay}, which holds for $d=3$ by
Remark~\ref{rM3:dim}, and $\varpi\le4$, $q_n>8$ for the absorption of
$e^{\varpi c/(e^2q_nn)}\le e^c$ into the constant.
\end{proof}

If the correlation length of the Villain model at large $\beta$ is of
the order of the inverse Debye mass, as the semiclassical picture
predicts, Proposition~\ref{rM3:xi} says that a constant $\varpi$ bounded
below is possible only on cubes with $n\gtrsim\varpi/(e^2q_nm_D)$, that is
of the size of the inverse photon mass. The lower bound on correlations
required in the proposition is not proved here for the Villain model. We
use from~\cite{GM82} only the area law and the exponential decay of
correlations, which are upper bounds, as summarised in~\cite{Cha18,RFN24};
whether~\cite{GM82} also gives a matching lower bound on plaquette
correlations, which would make the hypothesis of
Proposition~\ref{rM3:xi} hold with $\xi$ of the order of the inverse
photon mass, is not checked here. In the Debye--H\"uckel approximation the
corresponding necessary condition is the sharper quadratic bound of
Theorem~\ref{rM3:DH}(ii), $\varpi\le\frac{45}{32}(n\mathsf m)^2$ on every
torus.

\begin{remark}[comparison with the free fields and the Maxwell field]
\label{rM3:table}
In units of the cube, $s=$ mass $\times$ side, the cube constant behaves
as follows ($d=3$ here, general $d$ in Sec.~\ref{ssec:fffree}, $d=4$ in
Sec.~\ref{rD:sec}):
\begin{itemize}[leftmargin=*]
\item massive scalar: $\varpi\ge(\frac56)^d\frac{s^2}{s^2+10d}$
(Theorem~\ref{thm:ffscalar}), and an upper bound vanishing with $s$
(Proposition~\ref{prop:ffupper});
\item Maxwell field: no lower bound; $\varpi\le O(n^2/L^2)$
(Theorem~\ref{rD:maxwell});
\item Debye--H\"uckel field: $(\frac56)^3\frac{s^2}{2s^2+60}\le\varpi\le
\min(1+\frac2n,\frac{45}{32}s^2)$ on every torus (Theorem~\ref{rM3:DH});
\item Villain, plaquette algebra: lower bound open; $\varpi\le e^2q_nn/\xi$
whenever correlations decay no faster than $e^{-d/\xi}$
(Proposition~\ref{rM3:xi});
\item Villain, the functions $c\sin F_p$ of one plaquette: the inequality
holds with $\frac18$ for $\beta\ge2$ (Proposition~\ref{rM3:plaq}).
\end{itemize}
The Maxwell field and the Debye--H\"uckel field have the same exact part;
the difference is the term $\mathsf m^{-2}\lVert dF\rVert^2$, a mass for the
coexact (monopole) part only, and it is enough to turn the diffusive
failure into a uniform constant. The Villain field strength has the same
exact part again (Proposition~\ref{rM3:lawF}), and its coexact part is the
lattice-valued field $M$ instead of a Gaussian.
\end{remark}

\subsubsection{What is missing for the Villain measure}\label{rM3:missing}

The proof of Theorem~\ref{rM3:DH}(i) has two steps: the reduction of the
Dirichlet form to linear functionals of a Gaussian field
(Theorem~\ref{thm:ffchaos}, Lemma~\ref{lem:fflinear}), and a stable
decomposition of the precision form (Lemma~\ref{lem:fflions}). For the
Villain measure, Theorem~\ref{rM3:dualform} and
Proposition~\ref{rM3:interior} reduce $\mathrm{BP}(n,\varpi)$ on the
plaquette algebra to an inequality for $P_\beta=\gamma_\beta*\nu_{\mathcal M}$,
$\nu_{\mathcal M}$ the law of $M$, whose block law given the exterior is
the product of the Gaussian $N(0,\beta^{-1}\Pi_x)$ on $D_x$ and the
discrete Gaussian $\nu_{a_0}\propto e^{-\beta\lVert a\rVert^2/2}$ on the coset
$a_0+\mathcal M_x$, $a_0=a_0(F^{\rm ext}_x)$. Neither step transfers, and
the missing estimates are the following.

\emph{(O1) The class of functions.} The inequality is required for
periodic functions of $F$ only, and $P_\beta$ is not Gaussian, so there is
no chaos decomposition and no reduction to linear functionals; the linear
functionals $\ip cF$ of Theorem~\ref{rM3:DH}(ii), (iii) are not
observables of the gauge theory. The long-wavelength functions that decide the inequality
are nonlinear superpositions such as $\sum_y\cos(2\pi y_1/L)\sin F_{(y,12)}$
(Sec.~\ref{rM3:layersec}). What is missing is an inequality of Poincar\'e
or Brascamp--Lieb type for the block dynamics of $P_\beta$ on periodic
functions, with the monopole component treated non-perturbatively. Even
the upper bound $\varpi\lesssim(nm_D)^2$ for the Villain model, the
analogue of Theorem~\ref{rM3:DH}(ii), (iii), is not proved here: the
proof of Theorem~\ref{rM3:DH} would transfer only with a lower bound on
the variance of a periodic superposition such as the one above, uniform
in $L$, and an upper bound on $\E\Var(\sum_tN_t\mid F^{\rm ext}_x)$ for the
slab monopole numbers $N_t$ of Lemma~\ref{rM3:layer}.

\emph{(O2) Screening inside a block, conditionally on the exterior.} In the
Gaussian case the conditional covariance of the block does not depend on
the exterior. For the Villain measure the interior exact part is exactly
Gaussian and exterior-independent (Proposition~\ref{rM3:interior}(ii)),
so all dependence on the exterior sits in the coset $a_0+\mathcal M_x$.
The relaxation of a slow mode needs the slab monopole numbers $N_t$,
functions of the discrete component, to fluctuate under $\nu_{a_0}$ for
typical $a_0$; in the Debye--H\"uckel field the corresponding Gaussian
charge has variance at most $\mathsf m^2\lVert\mathbf 1_{T_t}\rVert^2/\beta$
(proof of Theorem~\ref{rM3:DH}(ii)). A form of the missing estimate
modelled on this order of magnitude, which we state as a target and do
not claim to be necessary or sufficient, is: with $\nu_{a_0}$ averaged over $a_0$ under $P_\beta$, a lower bound on the
block fluctuations of the discrete component of the form
$\E\Var(\ip{\mathsf n}{\tau}\mid F^{\rm ext}_x)\ge c\,\mathsf m^2\lVert\tau
\rVert^2/(4\pi^2\beta)$ for the profiles $\tau$ supported in the dual block,
with $\mathsf m\asymp m_D$ and $c$ independent of $L$, together with the
corresponding statement for nonlinear functions. On cubes of side $n\ll
1/m_D$ the Debye--H\"uckel field has $\varpi_{\rm DH}\le\frac{45}{32}
(n\mathsf m)^2$ (Theorem~\ref{rM3:DH}(ii)); no such upper bound is proved
for the Villain measure.

\emph{(O3) The relation to the known results.} Theorem~\ref{thm:avgdecay}
shows that $\mathrm{BP}(n,\varpi)$ implies exponential decay of
correlations; G\"opfert and Mack prove confinement and exponential decay
of correlations of the Villain model at every $\beta$~\cite{GM82}, as
summarised in~\cite{Cha18,RFN24}. The converse implication, from decay of
correlations of the measure to a spectral gap of the cube dynamics, is not
available here: the derivations of spectral gaps of block dynamics from
spatial mixing that we are aware of use mixing estimates that hold uniformly
in the boundary condition, and the decay statements quoted above concern
the infinite-volume state. The text of~\cite{GM82} could not be consulted
for this subsection, so whether its renormalisation-group estimates
contain bounds uniform in the boundary condition, which would be the input
for (O2), is not checked. Gaussian (sine-Gordon type) representations of
the kind used there represent expectations and not the cube dynamics
(Remark~\ref{rM3:nodual}).

\begin{remark}[the charged sectors]\label{rM3:charged}
By Lemma~\ref{rM3:split}, $\mathrm{BP}(n,\varpi)$ on $\mathcal A$ also needs
the inequality on $\mathcal A_F^\perp\cap\mathcal A$, which contains the
sectors $\mathcal A_\kappa$, $\kappa\ne0$, of functions transforming under
the shifts $T_c$ of Sec.~\ref{rM3:sectors}, such as the Polyakov loops
$e^{i\sum_{j=0}^{L-1}\theta_{(y+je_1,1)}}$, which lie in $\mathcal
A_{Le_1}$ (the loop has $L$ links of direction $1$). These
functions have mean zero, and a cube move changes the phase of a Polyakov
loop only through the links of the loop inside the cube. On a fixed torus
the Polyakov loops of $\SU(2)$ give the failure bound of
Theorem~\ref{thm:bpfixed}(ii); its $\U(1)$ analogue, and the behaviour of
these sectors on large tori, where confinement makes distant Polyakov
loops nearly independent, are not treated here. This subsection does not
claim $\mathcal A_0=\mathcal A_F$ either (Sec.~\ref{rM3:sectors}).
\end{remark}

\subsubsection{Status}\label{rM3:status}

Proved: the exact decomposition $F=G+M$ of the Villain field strength into
an independent massless Gaussian and a lattice-valued monopole field
(Proposition~\ref{rM3:lawF}); the dual form of $\mathrm{BP}(n,\varpi)$ on
the plaquette algebra (Theorem~\ref{rM3:dualform}); the block law given the
exterior (Proposition~\ref{rM3:interior}) and the lower bound for linear
functionals (Corollary~\ref{rM3:photon}); the exact layer-flux identity
with monopoles (Lemma~\ref{rM3:layer}); the inequality of $\mathrm{BP}$
with constant $\ge\frac18$ for the odd single-plaquette functions at $\beta
\ge2$ (Proposition~\ref{rM3:plaq}); for the Debye--H\"uckel Gaussian field of
mass $\mathsf m$, $\mathrm{BP}(n,\varpi)$ with $\varpi\ge(5/6)^3s^2/(2s^2+60)$
uniformly in $L$ and $\beta$, and $\varpi\le\min(1+2/n,\frac{45}{32}s^2)$,
$s=n\mathsf m$ (Theorem~\ref{rM3:DH}); and the necessary condition of
Proposition~\ref{rM3:xi}. Open: $\mathrm{BP}(n,\varpi)$ for the Villain
measure on any class of functions containing the long-wavelength
superpositions, on the plaquette algebra or on the charged sectors (items
(O1)--(O3) and Remark~\ref{rM3:charged}); also open here are the
fourth-moment bound that would give a $\beta$-independent constant for
$\cos F_p$ (Remark~\ref{rM3:plaqremark}), the identity $\mathcal A_0=
\mathcal A_F$ on the torus, and the lower bound on correlations in the
hypothesis of Proposition~\ref{rM3:xi}. Nothing in this subsection proves
$\mathrm{BP}(n,\varpi)$ or a mass gap for the three-dimensional Villain
model beyond what is quoted from~\cite{GM82}, nor a mass gap for the
four-dimensional theory, step (D), step (E), or any part of the Clay
problem; the three-dimensional abelian model serves only as the test case
in which the mechanism of Remark~\ref{rem:rSmissing} can be examined, and
in its Gaussian approximation the mechanism works.
